\documentclass[../../main.tex]{subfiles}

\begin{document}

\section{分离性质}\label{section:ch11-分离性质}

为了对拓扑空间以及这些空间之间的连续映射建立有趣的结果, 有必要丰富原始的拓扑结构. 本节我们对集合 $X$ 上的拓扑考虑分离性质, 它确保拓扑区分某种不交的点集, 作为一个结果, 确保 $X$ 上存在鲁棒的连续实值函数族.

在拓扑空间中我们已定义了点的邻域. 对于拓扑空间 $X$ 的子集 $K$, 谈到 $K$ 的邻域我们指的是包含 $K$ 的开集. 我们称 $X$ 的两个子集 $A$ 和 $B$ 被不交邻域分离, 若存在 $A$ 和 $B$ 的邻域, 它们是不交的. 对于拓扑空间 $X$ 我们考虑以下四种分离性质:

\begin{definition}[Tychonoff、Hausdorff、正则、正规分离性质]\label{definition:ch11-分离性质}
\begin{description}
\item[Tychonoff 分离性质] 对 $X$ 的两个点 $u$ 和 $v$, 存在 $u$ 的邻域不包含点 $v$ 且存在 $v$ 的邻域不包含点 $u$.
\item[Hausdorff 分离性质] $X$ 中每两个点可被不交的邻域分离.
\item[正则分离性质] Tychonoff 分离性质成立, 且进一步地, 每个闭集与不属于该集的点可被不交的邻域分离.
\item[正规分离性质] Tychonoff 分离性质成立, 且进一步地, 每两个不交的闭集可被不交的邻域分离.
\end{description}
\end{definition}

自然地, 我们称拓扑空间为 Tychonoff、Hausdorff、正则或正规, 若它满足相应的分离性质.

\begin{proposition}\label{proposition:ch11-Tychonoff等价}
拓扑空间 $X$ 是一个 Tychonoff 空间当且仅当每个单点集是闭的.
\end{proposition}

\begin{proof}
令 $x$ 为 $X$ 中的点. 集合 $\{x\}$ 是闭的当且仅当 $X\sim\{x\}$ 是开的. 现在 $X\sim\{x\}$ 是开的当且仅当对 $X\sim\{x\}$ 中的每个点 $y$, 存在 $y$ 的邻域不包含于 $X\sim\{x\}$, 即存在 $y$ 的邻域不包含 $x$.
\end{proof}

\begin{proposition}\label{proposition:ch11-度量空间正规}
每个度量空间是正规的.
\end{proposition}

\begin{proof}
令 $(X,\rho)$ 为度量空间. 定义 $X$ 的子集 $F$ 和 $X$ 中的点 $x$ 的距离为
\begin{align*}
\operatorname{dist}(x,F)=\inf\{\rho(x,x')\mid x'\text{ 属于 }F\}.
\end{align*}
令 $F_{1}$ 和 $F_{2}$ 为 $X$ 的闭的不交子集. 定义
\begin{align*}
\mathcal{O}_{1}=\{x\in X\mid\operatorname{dist}(x,F_{1})<\operatorname{dist}(x,F_{2})\}
\end{align*}
和
\begin{align*}
\mathcal{O}_{2}=\{x\in X\mid\operatorname{dist}(x,F_{2})<\operatorname{dist}(x,F_{1})\}.
\end{align*}
由于闭集的补集是开的, 若 $F$ 是闭的而 $x$ 不属于 $F$, 则 $\operatorname{dist}(x,F)>0$. 因此 $F_{1}\subseteq\mathcal{O}_{1}$, $F_{2}\subseteq\mathcal{O}_{2}$, 且显然 $\mathcal{O}_{1}\cap\mathcal{O}_{2}=\varnothing$. 此外, 用关于度量的三角不等式, 不难看出 $\mathcal{O}_{1}$ 与 $\mathcal{O}_{2}$ 是开的.
\end{proof}

利用显然的记号, 前两个命题给出了集合 $X$ 上的拓扑族之间的包含关系链:
\begin{align*}
\mathcal{T}_{\text{度量}}\subseteq\mathcal{T}_{\text{正规}}\subseteq\mathcal{T}_{\text{正则}}\subseteq\mathcal{T}_{\text{Hausdorff}}\subseteq\mathcal{T}_{\text{Tychonoff}}.
\end{align*}

我们用闭集的网邻域重述正规性结束这简短的一节.

\begin{proposition}\label{proposition:ch11-正规网邻域}
令 $X$ 为 Tychonoff 拓扑空间. 则 $X$ 是正规的当且仅当只要 $U$ 是 $X$ 的闭子集 $F$ 的邻域, 则存在 $F$ 的另一个邻域其闭包包含于 $U$, 即存在开集 $\mathcal{O}$ 使得
\begin{align*}
F\subseteq\mathcal{O}\subseteq\overline{\mathcal{O}}\subseteq U.
\end{align*}
\end{proposition}

\begin{proof}
首先假设 $X$ 是正规的. 由于 $F$ 和 $X\sim U$ 是不交的闭集, 存在不交的开集 $\mathcal{O}$ 和 $\mathcal{V}$ 使得 $F\subseteq\mathcal{O}$ 且 $X\sim U\subseteq\mathcal{V}$. 因此 $\mathcal{O}\subseteq X\sim\mathcal{V}\subseteq U$. 因为 $\mathcal{O}\subseteq X\sim\mathcal{V}$ 且 $X\sim\mathcal{V}$ 是闭的, 所以 $\overline{\mathcal{O}}\subseteq X\sim\mathcal{V}\subseteq U$.

为证明逆命题, 假定网邻域性质成立. 令 $A$ 和 $B$ 为 $X$ 的不交的闭子集. 则 $A\subseteq X\sim B$ 且 $X\sim B$ 是开的. 因此存在一个开集 $\mathcal{O}$ 使得 $A\subseteq\mathcal{O}\subseteq\overline{\mathcal{O}}\subseteq X\sim B$. 因此 $\mathcal{O}$ 和 $X\sim\overline{\mathcal{O}}$ 分别是 $A$ 和 $B$ 的不交邻域.
\end{proof}

\begin{exercise}\label{exercise:ch11-12}
证明: 若 $F$ 是正规空间 $X$ 的闭子集, 则子空间 $F$ 是正规的. $F$ 闭的假设是否必要?
\end{exercise}

\begin{exercise}\label{exercise:ch11-13}
令 $X$ 为拓扑空间. 证明 $X$ 是 Hausdorff 的当且仅当对角线 $D=\{(x_{1},x_{2})\in X\times X\mid x_{1}=x_{2}\}$ 是 $X\times X$ 的闭子集.
\end{exercise}

\begin{exercise}\label{exercise:ch11-14}
考虑具有由空集与形如 $(-\infty,c)$, $c\in\mathbf{R}$ 的集合组成的拓扑的实数集. 证明该空间是 Tychonoff 的但不是 Hausdorff 的.
\end{exercise}

\begin{exercise}\label{exercise:ch11-15}
(Zariski 拓扑)在 $\mathbf{R}^{n}$ 中令 $\mathcal{B}$ 为集族 $\{x\in\mathbf{R}^{n}\mid p(x)\neq 0\}$, 其中 $p$ 是 $n$ 个变量的多项式. 令 $\mathcal{T}$ 为 $X$ 上以 $\mathcal{B}$ 为子基的拓扑. 证明 $\mathcal{T}$ 是 $\mathbf{R}^{n}$ 的拓扑. 它是 Tychonoff 的但不是 Hausdorff 的.
\end{exercise}

\begin{exercise}\label{exercise:ch11-16}
证明 Sorgenfrey 直线与 Moore 平面是 Hausdorff 的(见习题 9 与 10).
\end{exercise}

\end{document}